% Options for packages loaded elsewhere
% Options for packages loaded elsewhere
\PassOptionsToPackage{unicode}{hyperref}
\PassOptionsToPackage{hyphens}{url}
%
\documentclass[
  ignorenonframetext,
  aspectratio=169,
  spanish,
]{beamer}
\newif\ifbibliography
\usepackage{pgfpages}
\setbeamertemplate{caption}[numbered]
\setbeamertemplate{caption label separator}{: }
\setbeamercolor{caption name}{fg=normal text.fg}
\beamertemplatenavigationsymbolshorizontal
% remove section numbering
\setbeamertemplate{part page}{
  \centering
  \begin{beamercolorbox}[sep=16pt,center]{part title}
    \usebeamerfont{part title}\insertpart\par
  \end{beamercolorbox}
}
\setbeamertemplate{section page}{
  \centering
  \begin{beamercolorbox}[sep=12pt,center]{section title}
    \usebeamerfont{section title}\insertsection\par
  \end{beamercolorbox}
}
\setbeamertemplate{subsection page}{
  \centering
  \begin{beamercolorbox}[sep=8pt,center]{subsection title}
    \usebeamerfont{subsection title}\insertsubsection\par
  \end{beamercolorbox}
}
% Prevent slide breaks in the middle of a paragraph
\widowpenalties 1 10000
\raggedbottom
\AtBeginPart{
  \frame{\partpage}
}
\AtBeginSection{
  \ifbibliography
  \else
    \frame{\sectionpage}
  \fi
}
\AtBeginSubsection{
  \frame{\subsectionpage}
}
\usepackage{iftex}
\ifPDFTeX
  \usepackage[T1]{fontenc}
  \usepackage[utf8]{inputenc}
  \usepackage{textcomp} % provide euro and other symbols
\else % if luatex or xetex
  \usepackage{unicode-math} % this also loads fontspec
  \defaultfontfeatures{Scale=MatchLowercase}
  \defaultfontfeatures[\rmfamily]{Ligatures=TeX,Scale=1}
\fi
\usepackage{lmodern}

\usetheme[]{default}
\usefonttheme[]{default}
\ifPDFTeX\else
  % xetex/luatex font selection
\fi
% Use upquote if available, for straight quotes in verbatim environments
\IfFileExists{upquote.sty}{\usepackage{upquote}}{}
\IfFileExists{microtype.sty}{% use microtype if available
  \usepackage[]{microtype}
  \UseMicrotypeSet[protrusion]{basicmath} % disable protrusion for tt fonts
}{}
\makeatletter
\@ifundefined{KOMAClassName}{% if non-KOMA class
  \IfFileExists{parskip.sty}{%
    \usepackage{parskip}
  }{% else
    \setlength{\parindent}{0pt}
    \setlength{\parskip}{6pt plus 2pt minus 1pt}}
}{% if KOMA class
  \KOMAoptions{parskip=half}}
\makeatother

\usepackage{color}
\usepackage{fancyvrb}
\newcommand{\VerbBar}{|}
\newcommand{\VERB}{\Verb[commandchars=\\\{\}]}
\DefineVerbatimEnvironment{Highlighting}{Verbatim}{commandchars=\\\{\}}
% Add ',fontsize=\small' for more characters per line
\usepackage{framed}
\definecolor{shadecolor}{RGB}{241,243,245}
\newenvironment{Shaded}{\begin{snugshade}}{\end{snugshade}}
\newcommand{\AlertTok}[1]{\textcolor[rgb]{0.68,0.00,0.00}{#1}}
\newcommand{\AnnotationTok}[1]{\textcolor[rgb]{0.37,0.37,0.37}{#1}}
\newcommand{\AttributeTok}[1]{\textcolor[rgb]{0.40,0.45,0.13}{#1}}
\newcommand{\BaseNTok}[1]{\textcolor[rgb]{0.68,0.00,0.00}{#1}}
\newcommand{\BuiltInTok}[1]{\textcolor[rgb]{0.00,0.23,0.31}{#1}}
\newcommand{\CharTok}[1]{\textcolor[rgb]{0.13,0.47,0.30}{#1}}
\newcommand{\CommentTok}[1]{\textcolor[rgb]{0.37,0.37,0.37}{#1}}
\newcommand{\CommentVarTok}[1]{\textcolor[rgb]{0.37,0.37,0.37}{\textit{#1}}}
\newcommand{\ConstantTok}[1]{\textcolor[rgb]{0.56,0.35,0.01}{#1}}
\newcommand{\ControlFlowTok}[1]{\textcolor[rgb]{0.00,0.23,0.31}{\textbf{#1}}}
\newcommand{\DataTypeTok}[1]{\textcolor[rgb]{0.68,0.00,0.00}{#1}}
\newcommand{\DecValTok}[1]{\textcolor[rgb]{0.68,0.00,0.00}{#1}}
\newcommand{\DocumentationTok}[1]{\textcolor[rgb]{0.37,0.37,0.37}{\textit{#1}}}
\newcommand{\ErrorTok}[1]{\textcolor[rgb]{0.68,0.00,0.00}{#1}}
\newcommand{\ExtensionTok}[1]{\textcolor[rgb]{0.00,0.23,0.31}{#1}}
\newcommand{\FloatTok}[1]{\textcolor[rgb]{0.68,0.00,0.00}{#1}}
\newcommand{\FunctionTok}[1]{\textcolor[rgb]{0.28,0.35,0.67}{#1}}
\newcommand{\ImportTok}[1]{\textcolor[rgb]{0.00,0.46,0.62}{#1}}
\newcommand{\InformationTok}[1]{\textcolor[rgb]{0.37,0.37,0.37}{#1}}
\newcommand{\KeywordTok}[1]{\textcolor[rgb]{0.00,0.23,0.31}{\textbf{#1}}}
\newcommand{\NormalTok}[1]{\textcolor[rgb]{0.00,0.23,0.31}{#1}}
\newcommand{\OperatorTok}[1]{\textcolor[rgb]{0.37,0.37,0.37}{#1}}
\newcommand{\OtherTok}[1]{\textcolor[rgb]{0.00,0.23,0.31}{#1}}
\newcommand{\PreprocessorTok}[1]{\textcolor[rgb]{0.68,0.00,0.00}{#1}}
\newcommand{\RegionMarkerTok}[1]{\textcolor[rgb]{0.00,0.23,0.31}{#1}}
\newcommand{\SpecialCharTok}[1]{\textcolor[rgb]{0.37,0.37,0.37}{#1}}
\newcommand{\SpecialStringTok}[1]{\textcolor[rgb]{0.13,0.47,0.30}{#1}}
\newcommand{\StringTok}[1]{\textcolor[rgb]{0.13,0.47,0.30}{#1}}
\newcommand{\VariableTok}[1]{\textcolor[rgb]{0.07,0.07,0.07}{#1}}
\newcommand{\VerbatimStringTok}[1]{\textcolor[rgb]{0.13,0.47,0.30}{#1}}
\newcommand{\WarningTok}[1]{\textcolor[rgb]{0.37,0.37,0.37}{\textit{#1}}}

\usepackage{longtable,booktabs,array}
\usepackage{calc} % for calculating minipage widths
\usepackage{caption}
% Make caption package work with longtable
\makeatletter
\def\fnum@table{\tablename~\thetable}
\makeatother
\usepackage{graphicx}
\makeatletter
\newsavebox\pandoc@box
\newcommand*\pandocbounded[1]{% scales image to fit in text height/width
  \sbox\pandoc@box{#1}%
  \Gscale@div\@tempa{\textheight}{\dimexpr\ht\pandoc@box+\dp\pandoc@box\relax}%
  \Gscale@div\@tempb{\linewidth}{\wd\pandoc@box}%
  \ifdim\@tempb\p@<\@tempa\p@\let\@tempa\@tempb\fi% select the smaller of both
  \ifdim\@tempa\p@<\p@\scalebox{\@tempa}{\usebox\pandoc@box}%
  \else\usebox{\pandoc@box}%
  \fi%
}
% Set default figure placement to htbp
\def\fps@figure{htbp}
\makeatother



\ifLuaTeX
\usepackage[bidi=basic]{babel}
\else
\usepackage[bidi=default]{babel}
\fi
% get rid of language-specific shorthands (see #6817):
\let\LanguageShortHands\languageshorthands
\def\languageshorthands#1{}


\setlength{\emergencystretch}{3em} % prevent overfull lines

\providecommand{\tightlist}{%
  \setlength{\itemsep}{0pt}\setlength{\parskip}{0pt}}



 


\setbeamertemplate{headline}{} % Elimina la cabecera del tema Madrid
\setbeamertemplate{section in head/foot}{} % Limpia la numeración de secciones
\renewcommand{\contentsname}{Índice}
\renewcommand{\refname}{Referencias}
\renewcommand{\figurename}{Figura}
\renewcommand{\tablename}{Tabla}
\logo{%
  \begin{tabular}{c}
    \includegraphics[height=1.2cm]{EscudoU.png}
  \end{tabular}%
}
% --- CONFIGURACIÓN DE COLORES INSTITUCIONALES ---
\definecolor{UnicorGreen}{HTML}{024426}
\definecolor{UnicorText}{HTML}{F9F9F9}
\definecolor{UnicorWhite}{HTML}{FFFFFF}
\definecolor{UnicorGold}{HTML}{D1B061}

% Aplicación al fondo y texto general
\setbeamercolor{background canvas}{bg=UnicorGreen}
\setbeamercolor{normal text}{fg=UnicorText}

% Aplicación a títulos y cabeceras
%\setbeamercolor{title}{fg=UnicorWhite}
\setbeamercolor{frametitle}{fg=UnicorWhite}
\setbeamercolor{framesubtitle}{fg=UnicorWhite}

% Aplicación a viñetas y estructura
\setbeamercolor{structure}{fg=UnicorWhite}
\setbeamercolor{itemize item}{fg=UnicorWhite}
\setbeamercolor{itemize subitem}{fg=UnicorWhite}
\setbeamercolor{enumerate item}{fg=UnicorWhite}

% Forzar a LaTeX a repintar el texto general al iniciar el documento
\AtBeginDocument{\usebeamercolor[fg]{normal text}}

% Aplicación a hipervínculos
\hypersetup{colorlinks=true, urlcolor=UnicorGold, linkcolor=UnicorWhite}

% --- ACTIVAR ESTILO DE BLOQUES Y SOMBRAS (ESTILO MADRID) ---
\useinnertheme[shadow=true]{rounded}

% --- COLORES DEL RECUADRO DEL TÍTULO PRINCIPAL ---
% Fondo dorado con texto verde
\setbeamercolor{title}{bg=UnicorGold, fg=UnicorGreen}
\setbeamercolor{frametitle}{fg=UnicorWhite} % Mantiene el título de las diapositivas normal
\setbeamercolor{framesubtitle}{fg=UnicorWhite}

% --- COLORES PARA TEOREMAS, DEFINICIONES Y BLOQUES ---
% Título del bloque (Fondo dorado, texto verde)
\setbeamercolor{block title}{bg=UnicorGold, fg=UnicorGreen}
% Cuerpo del bloque (Fondo blanco, texto verde)
\setbeamercolor{block body}{bg=UnicorWhite, fg=UnicorGreen}

\AtBeginDocument{
  \author{%
    \begin{tabular}{c}
      Mario A. Morales Rivera \\
      \footnotesize \texttt{mamorales@correo.unicordoba.edu.co} \\
      \footnotesize Universidad de Córdoba
    \end{tabular}
  }
  \institute{}
}
% --- LIMPIAR DIAPOSITIVAS DE SECCIÓN (ELIMINAR "SECTION X") ---
\setbeamertemplate{section page}{
  \begin{center}
    \begin{beamercolorbox}[sep=12pt,center,rounded=true,shadow=true]{title}
      \usebeamerfont{title}\insertsection\par
    \end{beamercolorbox}
  \end{center}
}
\makeatletter
\@ifpackageloaded{tcolorbox}{}{\usepackage[skins,breakable]{tcolorbox}}
\@ifpackageloaded{fontawesome5}{}{\usepackage{fontawesome5}}
\definecolor{quarto-callout-color}{HTML}{909090}
\definecolor{quarto-callout-note-color}{HTML}{0758E5}
\definecolor{quarto-callout-important-color}{HTML}{CC1914}
\definecolor{quarto-callout-warning-color}{HTML}{EB9113}
\definecolor{quarto-callout-tip-color}{HTML}{00A047}
\definecolor{quarto-callout-caution-color}{HTML}{FC5300}
\definecolor{quarto-callout-color-frame}{HTML}{acacac}
\definecolor{quarto-callout-note-color-frame}{HTML}{4582ec}
\definecolor{quarto-callout-important-color-frame}{HTML}{d9534f}
\definecolor{quarto-callout-warning-color-frame}{HTML}{f0ad4e}
\definecolor{quarto-callout-tip-color-frame}{HTML}{02b875}
\definecolor{quarto-callout-caution-color-frame}{HTML}{fd7e14}
\makeatother
\makeatletter
\@ifpackageloaded{caption}{}{\usepackage{caption}}
\AtBeginDocument{%
\ifdefined\contentsname
  \renewcommand*\contentsname{Tabla de contenidos}
\else
  \newcommand\contentsname{Tabla de contenidos}
\fi
\ifdefined\listfigurename
  \renewcommand*\listfigurename{Listado de Figuras}
\else
  \newcommand\listfigurename{Listado de Figuras}
\fi
\ifdefined\listtablename
  \renewcommand*\listtablename{Listado de Tablas}
\else
  \newcommand\listtablename{Listado de Tablas}
\fi
\ifdefined\figurename
  \renewcommand*\figurename{Figura}
\else
  \newcommand\figurename{Figura}
\fi
\ifdefined\tablename
  \renewcommand*\tablename{Tabla}
\else
  \newcommand\tablename{Tabla}
\fi
}
\@ifpackageloaded{float}{}{\usepackage{float}}
\floatstyle{ruled}
\@ifundefined{c@chapter}{\newfloat{codelisting}{h}{lop}}{\newfloat{codelisting}{h}{lop}[chapter]}
\floatname{codelisting}{Listado}
\newcommand*\listoflistings{\listof{codelisting}{Listado de Listados}}
\makeatother
\makeatletter
\makeatother
\makeatletter
\@ifpackageloaded{caption}{}{\usepackage{caption}}
\@ifpackageloaded{subcaption}{}{\usepackage{subcaption}}
\makeatother

\usepackage{bookmark}
\IfFileExists{xurl.sty}{\usepackage{xurl}}{} % add URL line breaks if available
\urlstyle{same}
\hypersetup{
  pdftitle={Data Frames},
  pdflang={es},
  hidelinks,
  pdfcreator={LaTeX via pandoc}}


\title{Data Frames}
\author{}
\date{}

\begin{document}
\frame{\titlepage}

\renewcommand*\contentsname{Tabla de contenidos}
\begin{frame}[allowframebreaks]
  \frametitle{Tabla de contenidos}
  \setcounter{tocdepth}{3}
  \tableofcontents
\end{frame}
\setcounter{tocdepth}{3}
\tableofcontents
}

\section{Data Frames conceptos
básicos}\label{data-frames-conceptos-buxe1sicos}

\begin{frame}{Data Frames en R}
\phantomsection\label{data-frames-en-r}
\begin{tcolorbox}[enhanced jigsaw, colframe=quarto-callout-note-color-frame, opacityback=0, breakable, rightrule=.15mm, leftrule=.75mm, coltitle=black, left=2mm, titlerule=0mm, toprule=.15mm, bottomtitle=1mm, colback=white, toptitle=1mm, title=\textcolor{quarto-callout-note-color}{\faInfo}\hspace{0.5em}{Nota}, arc=.35mm, colbacktitle=quarto-callout-note-color!10!white, bottomrule=.15mm, opacitybacktitle=0.6]

Un \textbf{Data Frame} en R es una estructura de datos bidimensional y
tabular, diseñada específicamente para el análisis estadístico.
Técnicamente, es una lista de vectores (o factores) que comparten
exactamente la misma longitud.

\end{tcolorbox}
\end{frame}

\begin{frame}{Data frames: características}
\phantomsection\label{data-frames-caracteruxedsticas}
\begin{itemize}
\item
  Cada \textbf{columna} representa una variable y cada \textbf{fila}
  representa una observación o caso.
\item
  A diferencia de las matrices en R (que admiten un solo tipo de dato),
  los data frames son heterogéneos: permiten almacenar diferentes tipos
  de datos (numéricos, caracteres, lógicos, factores) en distintas
  columnas.
\item
  Es el objeto estándar para la manipulación de datos y el ajuste de
  modelos estadísticos en R.
\end{itemize}
\end{frame}

\section{DataFrames: creación}\label{dataframes-creaciuxf3n}

\begin{frame}[fragile]{Creación de DataFrames}
\phantomsection\label{creaciuxf3n-de-dataframes}
Una forma de crear \texttt{DataFrames} en R es mediante la función
\texttt{data.frame}, como se muestra a continuación

\begin{Shaded}
\begin{Highlighting}[]
\FunctionTok{set.seed}\NormalTok{(}\DecValTok{1234}\NormalTok{)}
\CommentTok{\# 10 datos simulados de una variable aleatoria binmial con n=20 y p=0.4}
\NormalTok{x}\OtherTok{\textless{}{-}}\FunctionTok{rbinom}\NormalTok{(}\DecValTok{10}\NormalTok{,}\DecValTok{20}\NormalTok{,}\FloatTok{0.4}\NormalTok{)  }
\CommentTok{\# vector de caracteres aleatorios  }
\NormalTok{y}\OtherTok{\textless{}{-}}\FunctionTok{sample}\NormalTok{(}\FunctionTok{c}\NormalTok{(}\StringTok{"F"}\NormalTok{,}\StringTok{"M"}\NormalTok{),}\DecValTok{10}\NormalTok{,}\AttributeTok{replace=}\ConstantTok{TRUE}\NormalTok{)}
\CommentTok{\# vector logico }
\NormalTok{z}\OtherTok{\textless{}{-}}\FunctionTok{runif}\NormalTok{(}\DecValTok{10}\NormalTok{)}\SpecialCharTok{\textgreater{}}\FloatTok{0.5}
\end{Highlighting}
\end{Shaded}
\end{frame}

\begin{frame}[fragile]{}
\phantomsection\label{section}
A continuación creamos un \texttt{DataFrame} con estos tres vectores

\begin{Shaded}
\begin{Highlighting}[]
\NormalTok{datos}\OtherTok{=}\FunctionTok{data.frame}\NormalTok{(x,y,z)}
\NormalTok{datos}
\end{Highlighting}
\end{Shaded}

\begin{verbatim}
    x y     z
1   5 M FALSE
2   9 M FALSE
3   9 M FALSE
4   9 F FALSE
5  10 M FALSE
6   9 M  TRUE
7   3 M  TRUE
8   6 F  TRUE
9   9 M  TRUE
10  8 M FALSE
\end{verbatim}
\end{frame}

\section{Inspección rápida y
estructura}\label{inspecciuxf3n-ruxe1pida-y-estructura}

\begin{frame}[fragile]{Estructura del data}
\phantomsection\label{estructura-del-data}
Para verificar la estructura interna, los tipos de datos de cada columna
y la longitud del data frame, usamos \texttt{str(datos)}

\begin{Shaded}
\begin{Highlighting}[]
\FunctionTok{str}\NormalTok{(datos)}
\end{Highlighting}
\end{Shaded}

\begin{verbatim}
'data.frame':   10 obs. of  3 variables:
 $ x: int  5 9 9 9 10 9 3 6 9 8
 $ y: chr  "M" "M" "M" "F" ...
 $ z: logi  FALSE FALSE FALSE FALSE FALSE TRUE ...
\end{verbatim}
\end{frame}

\begin{frame}[fragile]{Resumen básico}
\phantomsection\label{resumen-buxe1sico}
Para obtener un resumen estadístico básico (mínimo, máximo, cuartiles,
media) para variables numéricas, y conteos para factores/caracteres.

\begin{Shaded}
\begin{Highlighting}[]
\FunctionTok{summary}\NormalTok{(datos) }
\end{Highlighting}
\end{Shaded}

\begin{verbatim}
       x             y                 z          
 Min.   : 3.0   Length:10          Mode :logical  
 1st Qu.: 6.5   Class :character   FALSE:6        
 Median : 9.0   Mode  :character   TRUE :4        
 Mean   : 7.7                                     
 3rd Qu.: 9.0                                     
 Max.   :10.0                                     
\end{verbatim}
\end{frame}

\begin{frame}[fragile]{Dimensiones del data}
\phantomsection\label{dimensiones-del-data}
Para saber rápidamente las dimensiones (filas y columnas) del data frame
usamos \texttt{dim(datos)}

\begin{Shaded}
\begin{Highlighting}[]
\FunctionTok{dim}\NormalTok{(datos)}
\end{Highlighting}
\end{Shaded}

\begin{verbatim}
[1] 10  3
\end{verbatim}
\end{frame}

\begin{frame}[fragile]{Las funciones \texttt{head()} y \texttt{tail()}}
\phantomsection\label{las-funciones-head-y-tail}
Para explorar las primeras o últimas filas de un data frame grande sin
saturar la consola.

\begin{Shaded}
\begin{Highlighting}[]
\FunctionTok{head}\NormalTok{(datos,}\DecValTok{3}\NormalTok{)}
\end{Highlighting}
\end{Shaded}

\begin{verbatim}
  x y     z
1 5 M FALSE
2 9 M FALSE
3 9 M FALSE
\end{verbatim}

\begin{Shaded}
\begin{Highlighting}[]
\FunctionTok{tail}\NormalTok{(datos)}
\end{Highlighting}
\end{Shaded}

\begin{verbatim}
    x y     z
5  10 M FALSE
6   9 M  TRUE
7   3 M  TRUE
8   6 F  TRUE
9   9 M  TRUE
10  8 M FALSE
\end{verbatim}
\end{frame}

\begin{frame}[fragile]{¿Es un data?}
\phantomsection\label{es-un-data}
Para saber si efectivamente un objeto en R es un \texttt{DataFrame}
podemos ejecutar los siguientes comandos

\begin{Shaded}
\begin{Highlighting}[]
\FunctionTok{class}\NormalTok{(datos) }
\end{Highlighting}
\end{Shaded}

\begin{verbatim}
[1] "data.frame"
\end{verbatim}

\begin{Shaded}
\begin{Highlighting}[]
\FunctionTok{is.data.frame}\NormalTok{(datos) }
\end{Highlighting}
\end{Shaded}

\begin{verbatim}
[1] TRUE
\end{verbatim}

\begin{Shaded}
\begin{Highlighting}[]
\FunctionTok{mode}\NormalTok{(datos) }
\end{Highlighting}
\end{Shaded}

\begin{verbatim}
[1] "list"
\end{verbatim}
\end{frame}

\section{Nombres de filas y columnas}\label{nombres-de-filas-y-columnas}

\begin{frame}[fragile]{Nombres de columnas}
\phantomsection\label{nombres-de-columnas}
Observe que los nombres de columna se toman de los nombres de los
vectores usados para crear el data, si queremos otros nombres de
columnas procedemos de la siguiente forma

\begin{Shaded}
\begin{Highlighting}[]
\NormalTok{datos}\OtherTok{=}\FunctionTok{data.frame}\NormalTok{(}\AttributeTok{Cx=}\NormalTok{x,}\AttributeTok{Cy=}\NormalTok{y,}\AttributeTok{Cz=}\NormalTok{z) }
\NormalTok{datos}
\end{Highlighting}
\end{Shaded}

\begin{verbatim}
   Cx Cy    Cz
1   5  M FALSE
2   9  M FALSE
3   9  M FALSE
4   9  F FALSE
5  10  M FALSE
6   9  M  TRUE
7   3  M  TRUE
8   6  F  TRUE
9   9  M  TRUE
10  8  M FALSE
\end{verbatim}
\end{frame}

\begin{frame}[fragile]{Nombres de filas}
\phantomsection\label{nombres-de-filas}
También podemos colocar nombres de filas al data al momento de crearla,
de la siguiente forma

\begin{Shaded}
\begin{Highlighting}[]
\NormalTok{datos}\OtherTok{=}\FunctionTok{data.frame}\NormalTok{(}\AttributeTok{Cx=}\NormalTok{x,}\AttributeTok{Cy=}\NormalTok{y,}\AttributeTok{Cz=}\NormalTok{z,}\AttributeTok{row.names=}\NormalTok{month.abb[}\DecValTok{1}\SpecialCharTok{:}\DecValTok{10}\NormalTok{])}
\NormalTok{datos  }
\end{Highlighting}
\end{Shaded}

\begin{verbatim}
    Cx Cy    Cz
Jan  5  M FALSE
Feb  9  M FALSE
Mar  9  M FALSE
Apr  9  F FALSE
May 10  M FALSE
Jun  9  M  TRUE
Jul  3  M  TRUE
Aug  6  F  TRUE
Sep  9  M  TRUE
Oct  8  M FALSE
\end{verbatim}
\end{frame}

\begin{frame}[fragile]{Nombres de filas II}
\phantomsection\label{nombres-de-filas-ii}
Cuando uno de los vectores con que creamos el data tiene nombres para
sus elementos, la función \texttt{data.frame} toma estos nombres como
nombres de fila del data creado.

\begin{Shaded}
\begin{Highlighting}[]
\FunctionTok{names}\NormalTok{(y)}\OtherTok{=}\NormalTok{LETTERS[}\DecValTok{1}\SpecialCharTok{:}\DecValTok{10}\NormalTok{] }
\NormalTok{datos}\OtherTok{=}\FunctionTok{data.frame}\NormalTok{(}\AttributeTok{Cx=}\NormalTok{x,}\AttributeTok{Cy=}\NormalTok{y,}\AttributeTok{Cz=}\NormalTok{z) }
\NormalTok{datos}
\end{Highlighting}
\end{Shaded}

\begin{verbatim}
  Cx Cy    Cz
A  5  M FALSE
B  9  M FALSE
C  9  M FALSE
D  9  F FALSE
E 10  M FALSE
F  9  M  TRUE
G  3  M  TRUE
H  6  F  TRUE
I  9  M  TRUE
J  8  M FALSE
\end{verbatim}
\end{frame}

\begin{frame}[fragile]{Nombres de filas II}
\phantomsection\label{nombres-de-filas-ii-1}
¿Qué pasa si más de un vector tiene nombres? Veamos

\begin{Shaded}
\begin{Highlighting}[]
\FunctionTok{names}\NormalTok{(y)}\OtherTok{=}\NormalTok{LETTERS[}\DecValTok{1}\SpecialCharTok{:}\DecValTok{10}\NormalTok{] }
\FunctionTok{names}\NormalTok{(x)}\OtherTok{=}\NormalTok{month.abb[}\DecValTok{1}\SpecialCharTok{:}\DecValTok{10}\NormalTok{]}
\NormalTok{datos}\OtherTok{=}\FunctionTok{data.frame}\NormalTok{(}\AttributeTok{Cy=}\NormalTok{y,}\AttributeTok{Cx=}\NormalTok{x,}\AttributeTok{Cz=}\NormalTok{z,}\AttributeTok{row.names=}\FunctionTok{names}\NormalTok{(x)) }
\NormalTok{datos}
\end{Highlighting}
\end{Shaded}

\begin{verbatim}
    Cy Cx    Cz
Jan  M  5 FALSE
Feb  M  9 FALSE
Mar  M  9 FALSE
Apr  F  9 FALSE
May  M 10 FALSE
Jun  M  9  TRUE
Jul  M  3  TRUE
Aug  F  6  TRUE
Sep  M  9  TRUE
Oct  M  8 FALSE
\end{verbatim}
\end{frame}

\begin{frame}[fragile]{Nombres de filas IV}
\phantomsection\label{nombres-de-filas-iv}
Si no quiero nombres de fila se lo indico de la siguiente forma

\begin{Shaded}
\begin{Highlighting}[]
\NormalTok{datos}\OtherTok{=}\FunctionTok{data.frame}\NormalTok{(}\AttributeTok{Cy=}\NormalTok{y,}\AttributeTok{Cx=}\NormalTok{x,}\AttributeTok{Cz=}\NormalTok{z,}\AttributeTok{row.names=}\ConstantTok{NULL}\NormalTok{) }
\NormalTok{datos}
\end{Highlighting}
\end{Shaded}

\begin{verbatim}
   Cy Cx    Cz
1   M  5 FALSE
2   M  9 FALSE
3   M  9 FALSE
4   F  9 FALSE
5   M 10 FALSE
6   M  9  TRUE
7   M  3  TRUE
8   F  6  TRUE
9   M  9  TRUE
10  M  8 FALSE
\end{verbatim}
\end{frame}

\begin{frame}[fragile]{Nombres de filas, usos}
\phantomsection\label{nombres-de-filas-usos}
Los nombres de filas se pueden usar para recuperar filas del data
(indexamos el data con un vector de caracteres que contiene los nombres
de las filas que quiere recuperar)

\begin{Shaded}
\begin{Highlighting}[]
\NormalTok{datos}\OtherTok{=}\FunctionTok{data.frame}\NormalTok{(}\AttributeTok{Cx=}\NormalTok{x,}\AttributeTok{Cy=}\NormalTok{y,}\AttributeTok{Cz=}\NormalTok{z)  }
\NormalTok{datos[}\FunctionTok{c}\NormalTok{(}\StringTok{\textquotesingle{}Jan\textquotesingle{}}\NormalTok{,}\StringTok{\textquotesingle{}Mar\textquotesingle{}}\NormalTok{,}\StringTok{\textquotesingle{}May\textquotesingle{}}\NormalTok{),] }
\end{Highlighting}
\end{Shaded}

\begin{verbatim}
    Cx Cy    Cz
Jan  5  M FALSE
Mar  9  M FALSE
May 10  M FALSE
\end{verbatim}
\end{frame}

\begin{frame}[fragile]{Recuperación de nombres de filas}
\phantomsection\label{recuperaciuxf3n-de-nombres-de-filas}
Para recuperar los nombres de filas de un data usamos la función
\texttt{row.names} o \texttt{rownames}

\begin{Shaded}
\begin{Highlighting}[]
\CommentTok{\# filas }
\FunctionTok{row.names}\NormalTok{(datos) }
\end{Highlighting}
\end{Shaded}

\begin{verbatim}
 [1] "Jan" "Feb" "Mar" "Apr" "May" "Jun" "Jul" "Aug" "Sep" "Oct"
\end{verbatim}

\begin{Shaded}
\begin{Highlighting}[]
\FunctionTok{rownames}\NormalTok{(datos)}
\end{Highlighting}
\end{Shaded}

\begin{verbatim}
 [1] "Jan" "Feb" "Mar" "Apr" "May" "Jun" "Jul" "Aug" "Sep" "Oct"
\end{verbatim}

\begin{Shaded}
\begin{Highlighting}[]
\CommentTok{\# columnas }
\FunctionTok{colnames}\NormalTok{(datos) }
\end{Highlighting}
\end{Shaded}

\begin{verbatim}
[1] "Cx" "Cy" "Cz"
\end{verbatim}
\end{frame}

\begin{frame}[fragile]{Recuperar nombres de columnas}
\phantomsection\label{recuperar-nombres-de-columnas}
Para recuperar los dombres de las columnas usamos \texttt{colnames}, de
la siguiente forma

\begin{Shaded}
\begin{Highlighting}[]
\CommentTok{\# columnas }
\FunctionTok{colnames}\NormalTok{(datos) }
\end{Highlighting}
\end{Shaded}

\begin{verbatim}
[1] "Cx" "Cy" "Cz"
\end{verbatim}
\end{frame}

\begin{frame}[fragile]{Nombres de filas y columnas al tiempo}
\phantomsection\label{nombres-de-filas-y-columnas-al-tiempo}
La función \texttt{dimnames()} regresa una lista cuya primera componente
tiene el vector de nombres de fila y la segunda el vector de nombres de
columnas.

\begin{Shaded}
\begin{Highlighting}[]
\FunctionTok{dimnames}\NormalTok{(datos)}
\end{Highlighting}
\end{Shaded}

\begin{verbatim}
[[1]]
 [1] "Jan" "Feb" "Mar" "Apr" "May" "Jun" "Jul" "Aug" "Sep" "Oct"

[[2]]
[1] "Cx" "Cy" "Cz"
\end{verbatim}
\end{frame}

\begin{frame}[fragile]{Cambio de nombre}
\phantomsection\label{cambio-de-nombre}
Para cambiar los nombres de columnas procedemos de la siguiente forma

\begin{Shaded}
\begin{Highlighting}[]
\FunctionTok{colnames}\NormalTok{(datos)}\OtherTok{=}\FunctionTok{c}\NormalTok{(}\StringTok{\textquotesingle{}Colx\textquotesingle{}}\NormalTok{,}\StringTok{\textquotesingle{}Coly\textquotesingle{}}\NormalTok{,}\StringTok{\textquotesingle{}Colz\textquotesingle{}}\NormalTok{) }
\NormalTok{datos}
\end{Highlighting}
\end{Shaded}

\begin{verbatim}
    Colx Coly  Colz
Jan    5    M FALSE
Feb    9    M FALSE
Mar    9    M FALSE
Apr    9    F FALSE
May   10    M FALSE
Jun    9    M  TRUE
Jul    3    M  TRUE
Aug    6    F  TRUE
Sep    9    M  TRUE
Oct    8    M FALSE
\end{verbatim}
\end{frame}

\begin{frame}[fragile]{Cambiar el nombre de una sola columna}
\phantomsection\label{cambiar-el-nombre-de-una-sola-columna}
\begin{Shaded}
\begin{Highlighting}[]
\FunctionTok{colnames}\NormalTok{(datos)[}\DecValTok{2}\NormalTok{]}\OtherTok{=}\StringTok{\textquotesingle{}Col2\textquotesingle{}} 
\NormalTok{datos}
\end{Highlighting}
\end{Shaded}

\begin{verbatim}
    Colx Col2  Colz
Jan    5    M FALSE
Feb    9    M FALSE
Mar    9    M FALSE
Apr    9    F FALSE
May   10    M FALSE
Jun    9    M  TRUE
Jul    3    M  TRUE
Aug    6    F  TRUE
Sep    9    M  TRUE
Oct    8    M FALSE
\end{verbatim}
\end{frame}

\section{Adición de columnas}\label{adiciuxf3n-de-columnas}

\begin{frame}[fragile]{Creación de nuevas columnas, uso de \texttt{\$}}
\phantomsection\label{creaciuxf3n-de-nuevas-columnas-uso-de}
Podemos crear nuevas columnas en un data frame, a partir de las columnas
existentes, de la siguiente forma

\begin{Shaded}
\begin{Highlighting}[]
\NormalTok{datos}\SpecialCharTok{$}\NormalTok{Cx2}\OtherTok{=}\NormalTok{datos}\SpecialCharTok{$}\NormalTok{Colx}\SpecialCharTok{\^{}}\DecValTok{2} 
\NormalTok{datos}
\end{Highlighting}
\end{Shaded}

\begin{verbatim}
    Colx Col2  Colz Cx2
Jan    5    M FALSE  25
Feb    9    M FALSE  81
Mar    9    M FALSE  81
Apr    9    F FALSE  81
May   10    M FALSE 100
Jun    9    M  TRUE  81
Jul    3    M  TRUE   9
Aug    6    F  TRUE  36
Sep    9    M  TRUE  81
Oct    8    M FALSE  64
\end{verbatim}
\end{frame}

\begin{frame}[fragile]{Uso de
\texttt{{[}\textquotesingle{}nombre\textquotesingle{}{]}}}
\phantomsection\label{uso-de-nombre}
Crearemos una columna que contiene la columna \texttt{Col2} codificada
de la siguiente forma: Si \texttt{Col2} es \texttt{M} asignamos 1 (uno),
si es \texttt{F} asignamos 0 (cero).

\begin{Shaded}
\begin{Highlighting}[]
\NormalTok{datos[}\StringTok{"Sex"}\NormalTok{]}\OtherTok{=}\FunctionTok{ifelse}\NormalTok{(datos}\SpecialCharTok{$}\NormalTok{Col2}\SpecialCharTok{==}\StringTok{"M"}\NormalTok{,}\DecValTok{1}\NormalTok{,}\DecValTok{0}\NormalTok{)}
\NormalTok{datos}
\end{Highlighting}
\end{Shaded}

\begin{verbatim}
    Colx Col2  Colz Cx2 Sex
Jan    5    M FALSE  25   1
Feb    9    M FALSE  81   1
Mar    9    M FALSE  81   1
Apr    9    F FALSE  81   0
May   10    M FALSE 100   1
Jun    9    M  TRUE  81   1
Jul    3    M  TRUE   9   1
Aug    6    F  TRUE  36   0
Sep    9    M  TRUE  81   1
Oct    8    M FALSE  64   1
\end{verbatim}
\end{frame}

\begin{frame}[fragile]{Borrado de columna}
\phantomsection\label{borrado-de-columna}
Para borrar una columna usamos \texttt{NULL} de la siguiente forma

\begin{Shaded}
\begin{Highlighting}[]
\NormalTok{datos}\SpecialCharTok{$}\NormalTok{Cx2}\OtherTok{=}\ConstantTok{NULL}
\NormalTok{datos }
\end{Highlighting}
\end{Shaded}

\begin{verbatim}
    Colx Col2  Colz Sex
Jan    5    M FALSE   1
Feb    9    M FALSE   1
Mar    9    M FALSE   1
Apr    9    F FALSE   0
May   10    M FALSE   1
Jun    9    M  TRUE   1
Jul    3    M  TRUE   1
Aug    6    F  TRUE   0
Sep    9    M  TRUE   1
Oct    8    M FALSE   1
\end{verbatim}
\end{frame}

\section{Indexado}\label{indexado}

\begin{frame}{Indexando Data Frames}
\phantomsection\label{indexando-data-frames}
Como con los vectores y matrices, podemos indexar un data frame usando
como indices

\begin{itemize}
\tightlist
\item
  Enteros positivos
\item
  Enteros negativos
\item
  Valores lógicos
\item
  Caracteres
\end{itemize}

Esto es válido tanto para filas como para columnas
\end{frame}

\begin{frame}[fragile]{Veamos algunos ejemplos}
\phantomsection\label{veamos-algunos-ejemplos}
Seleccionamos las filas con índice 1, 5, 7, 9, 10 y las columnas con
índice 2 y 3

\begin{Shaded}
\begin{Highlighting}[]
\NormalTok{datos[}\FunctionTok{c}\NormalTok{(}\DecValTok{1}\NormalTok{,}\DecValTok{5}\NormalTok{,}\DecValTok{7}\NormalTok{,}\DecValTok{9}\NormalTok{,}\DecValTok{10}\NormalTok{),}\FunctionTok{c}\NormalTok{(}\DecValTok{2}\NormalTok{,}\DecValTok{3}\NormalTok{)] }
\end{Highlighting}
\end{Shaded}

\begin{verbatim}
    Col2  Colz
Jan    M FALSE
May    M FALSE
Jul    M  TRUE
Sep    M  TRUE
Oct    M FALSE
\end{verbatim}
\end{frame}

\begin{frame}[fragile]{Indexado con números negativos}
\phantomsection\label{indexado-con-nuxfameros-negativos}
Seleccionamos las filas con índice \texttt{distintos} a 1, 5, 7, 9, 10 y
las columnas con índice \texttt{distintos} a 2 y 3

\begin{Shaded}
\begin{Highlighting}[]
\NormalTok{datos[}\SpecialCharTok{{-}}\FunctionTok{c}\NormalTok{(}\DecValTok{1}\NormalTok{,}\DecValTok{5}\NormalTok{,}\DecValTok{7}\NormalTok{,}\DecValTok{9}\NormalTok{,}\DecValTok{10}\NormalTok{),}\SpecialCharTok{{-}}\FunctionTok{c}\NormalTok{(}\DecValTok{2}\NormalTok{,}\DecValTok{3}\NormalTok{)]}
\end{Highlighting}
\end{Shaded}

\begin{verbatim}
    Colx Sex
Feb    9   1
Mar    9   1
Apr    9   0
Jun    9   1
Aug    6   0
\end{verbatim}
\end{frame}

\begin{frame}[fragile]{Ejemplo de indexado con caracteres}
\phantomsection\label{ejemplo-de-indexado-con-caracteres}
Válido cuando las filas tienen nombres

\begin{Shaded}
\begin{Highlighting}[]
\NormalTok{datos[}\FunctionTok{c}\NormalTok{(}\StringTok{\textquotesingle{}Feb\textquotesingle{}}\NormalTok{,}\StringTok{\textquotesingle{}Apr\textquotesingle{}}\NormalTok{,}\StringTok{\textquotesingle{}Jun\textquotesingle{}}\NormalTok{,}\StringTok{\textquotesingle{}Aug\textquotesingle{}}\NormalTok{),}\SpecialCharTok{{-}}\DecValTok{1}\NormalTok{]  }
\end{Highlighting}
\end{Shaded}

\begin{verbatim}
    Col2  Colz Sex
Feb    M FALSE   1
Apr    F FALSE   0
Jun    M  TRUE   1
Aug    F  TRUE   0
\end{verbatim}
\end{frame}

\begin{frame}[fragile]{Indexación lógica}
\phantomsection\label{indexaciuxf3n-luxf3gica}
En este tipo de indexación se extraen las filas o columnas que coincidan
con un valor \texttt{TRUE} , veamos ejemplos

\begin{Shaded}
\begin{Highlighting}[]
\NormalTok{datos[}\FunctionTok{c}\NormalTok{(}\ConstantTok{TRUE}\NormalTok{,}\ConstantTok{FALSE}\NormalTok{,}\ConstantTok{TRUE}\NormalTok{,}\ConstantTok{FALSE}\NormalTok{,}\ConstantTok{TRUE}\NormalTok{,}\ConstantTok{FALSE}\NormalTok{,}\ConstantTok{TRUE}\NormalTok{,}\ConstantTok{FALSE}\NormalTok{,}\ConstantTok{TRUE}\NormalTok{,}\ConstantTok{FALSE}\NormalTok{), ]  }
\end{Highlighting}
\end{Shaded}

\begin{verbatim}
    Colx Col2  Colz Sex
Jan    5    M FALSE   1
Mar    9    M FALSE   1
May   10    M FALSE   1
Jul    3    M  TRUE   1
Sep    9    M  TRUE   1
\end{verbatim}

\begin{Shaded}
\begin{Highlighting}[]
\NormalTok{datos[}\SpecialCharTok{!}\FunctionTok{c}\NormalTok{(}\ConstantTok{TRUE}\NormalTok{,}\ConstantTok{FALSE}\NormalTok{,}\ConstantTok{TRUE}\NormalTok{,}\ConstantTok{FALSE}\NormalTok{,}\ConstantTok{TRUE}\NormalTok{,}\ConstantTok{FALSE}\NormalTok{,}\ConstantTok{TRUE}\NormalTok{,}\ConstantTok{FALSE}\NormalTok{,}\ConstantTok{TRUE}\NormalTok{,}\ConstantTok{FALSE}\NormalTok{), ]  }
\end{Highlighting}
\end{Shaded}

\begin{verbatim}
    Colx Col2  Colz Sex
Feb    9    M FALSE   1
Apr    9    F FALSE   0
Jun    9    M  TRUE   1
Aug    6    F  TRUE   0
Oct    8    M FALSE   1
\end{verbatim}
\end{frame}

\begin{frame}[fragile]{Ejercicio}
\phantomsection\label{ejercicio}
Ejercicio: Calcular la media de \texttt{Colx} pero para aquellas filas
conde \texttt{Col2} es igual a ``M'' y separadamente para donde
\texttt{Col2} es igual a ``F''

\begin{Shaded}
\begin{Highlighting}[]
\NormalTok{datosM}\OtherTok{=}\NormalTok{datos[datos}\SpecialCharTok{$}\NormalTok{Col2}\SpecialCharTok{==}\StringTok{"M"}\NormalTok{,] }
\FunctionTok{mean}\NormalTok{(datosM}\SpecialCharTok{$}\NormalTok{Colx) }
\end{Highlighting}
\end{Shaded}

\begin{verbatim}
[1] 7.75
\end{verbatim}

Para los femeninos, es decir, donde \texttt{Col2} es igual a ``F''

\begin{Shaded}
\begin{Highlighting}[]
\NormalTok{datosF}\OtherTok{=}\NormalTok{datos[datos}\SpecialCharTok{$}\NormalTok{Col2}\SpecialCharTok{==}\StringTok{"F"}\NormalTok{,]  }
\FunctionTok{mean}\NormalTok{(datosF}\SpecialCharTok{$}\NormalTok{Colx) }
\end{Highlighting}
\end{Shaded}

\begin{verbatim}
[1] 7.5
\end{verbatim}
\end{frame}

\begin{frame}[fragile]{Otra forma, usando la negación}
\phantomsection\label{otra-forma-usando-la-negaciuxf3n}
\begin{Shaded}
\begin{Highlighting}[]
\NormalTok{datosF}\OtherTok{=}\NormalTok{datos[}\SpecialCharTok{!}\NormalTok{datos}\SpecialCharTok{$}\NormalTok{Col2}\SpecialCharTok{==}\StringTok{"M"}\NormalTok{,]  }
\FunctionTok{mean}\NormalTok{(datosF}\SpecialCharTok{$}\NormalTok{Colx) }
\end{Highlighting}
\end{Shaded}

\begin{verbatim}
[1] 7.5
\end{verbatim}
\end{frame}

\begin{frame}[fragile]{Ejercicio}
\phantomsection\label{ejercicio-1}
Calcular la varianza de \texttt{Colx} para aquellas filas donde
\texttt{Colz} es TRUE

\begin{Shaded}
\begin{Highlighting}[]
\FunctionTok{var}\NormalTok{(datos[datos}\SpecialCharTok{$}\NormalTok{Colz,][[}\StringTok{"Colx"}\NormalTok{]]) }
\end{Highlighting}
\end{Shaded}

\begin{verbatim}
[1] 8.25
\end{verbatim}
\end{frame}

\begin{frame}[fragile]{Ejercicio}
\phantomsection\label{ejercicio-2}
Calcular la varianza de \texttt{Colx} para aquellas filas donde
\texttt{Colz} es FALSE

\begin{Shaded}
\begin{Highlighting}[]
\FunctionTok{var}\NormalTok{(datos[}\SpecialCharTok{!}\NormalTok{datos}\SpecialCharTok{$}\NormalTok{Colz,][[}\StringTok{"Colx"}\NormalTok{]])  }
\end{Highlighting}
\end{Shaded}

\begin{verbatim}
[1] 3.066667
\end{verbatim}
\end{frame}

\begin{frame}[fragile]{Ejercicio}
\phantomsection\label{ejercicio-3}
Calcular la media de \texttt{Colx} para equellas filas donde
\texttt{Col2} es `M' y \texttt{Colz} es FALSE

\begin{Shaded}
\begin{Highlighting}[]
\FunctionTok{mean}\NormalTok{(datos[datos}\SpecialCharTok{$}\NormalTok{Col2}\SpecialCharTok{==}\StringTok{\textquotesingle{}M\textquotesingle{}} \SpecialCharTok{\&} \SpecialCharTok{!}\NormalTok{datos}\SpecialCharTok{$}\NormalTok{Colz, ][[}\StringTok{\textquotesingle{}Colx\textquotesingle{}}\NormalTok{]])}
\end{Highlighting}
\end{Shaded}

\begin{verbatim}
[1] 8.2
\end{verbatim}
\end{frame}

\begin{frame}[fragile]{Ejercicio}
\phantomsection\label{ejercicio-4}
Cuántos hombres, es decir Col2==`M' tienen un valor de \texttt{Colx} por
encima de la media de \texttt{Colx}

\begin{Shaded}
\begin{Highlighting}[]
\FunctionTok{nrow}\NormalTok{(datos[datos}\SpecialCharTok{$}\NormalTok{Colx}\SpecialCharTok{\textgreater{}}\FunctionTok{mean}\NormalTok{(datos}\SpecialCharTok{$}\NormalTok{Colx) }\SpecialCharTok{\&}\NormalTok{ datos}\SpecialCharTok{$}\NormalTok{Col2}\SpecialCharTok{==}\StringTok{"M"}\NormalTok{ ,]  ) }
\end{Highlighting}
\end{Shaded}

\begin{verbatim}
[1] 6
\end{verbatim}
\end{frame}

\begin{frame}[fragile]{otra forma}
\phantomsection\label{otra-forma}
\begin{Shaded}
\begin{Highlighting}[]
\FunctionTok{sum}\NormalTok{(datos}\SpecialCharTok{$}\NormalTok{Colx}\SpecialCharTok{\textgreater{}}\FunctionTok{mean}\NormalTok{(datos}\SpecialCharTok{$}\NormalTok{Colx) }\SpecialCharTok{\&}\NormalTok{ datos}\SpecialCharTok{$}\NormalTok{Col2}\SpecialCharTok{==}\StringTok{"M"}\NormalTok{)}
\end{Highlighting}
\end{Shaded}

\begin{verbatim}
[1] 6
\end{verbatim}
\end{frame}

\begin{frame}[fragile]{otra forma}
\phantomsection\label{otra-forma-1}
\begin{Shaded}
\begin{Highlighting}[]
\FunctionTok{sum}\NormalTok{(datos[datos}\SpecialCharTok{$}\NormalTok{Colx}\SpecialCharTok{\textgreater{}}\FunctionTok{mean}\NormalTok{(datos}\SpecialCharTok{$}\NormalTok{Colx),][[}\StringTok{\textquotesingle{}Col2\textquotesingle{}}\NormalTok{]]}\SpecialCharTok{==}\StringTok{"M"}\NormalTok{  )}
\end{Highlighting}
\end{Shaded}

\begin{verbatim}
[1] 6
\end{verbatim}
\end{frame}




\end{document}
